% Options for packages loaded elsewhere
% Options for packages loaded elsewhere
\PassOptionsToPackage{unicode}{hyperref}
\PassOptionsToPackage{hyphens}{url}
\PassOptionsToPackage{dvipsnames,svgnames,x11names}{xcolor}
%
\documentclass[
  polish,
  letterpaper,
  DIV=11,
  numbers=noendperiod]{scrartcl}
\usepackage{xcolor}
\usepackage{amsmath,amssymb}
\setcounter{secnumdepth}{-\maxdimen} % remove section numbering
\usepackage{iftex}
\ifPDFTeX
  \usepackage[T1]{fontenc}
  \usepackage[utf8]{inputenc}
  \usepackage{textcomp} % provide euro and other symbols
\else % if luatex or xetex
  \usepackage{unicode-math} % this also loads fontspec
  \defaultfontfeatures{Scale=MatchLowercase}
  \defaultfontfeatures[\rmfamily]{Ligatures=TeX,Scale=1}
\fi
\usepackage{lmodern}
\ifPDFTeX\else
  % xetex/luatex font selection
\fi
% Use upquote if available, for straight quotes in verbatim environments
\IfFileExists{upquote.sty}{\usepackage{upquote}}{}
\IfFileExists{microtype.sty}{% use microtype if available
  \usepackage[]{microtype}
  \UseMicrotypeSet[protrusion]{basicmath} % disable protrusion for tt fonts
}{}
\makeatletter
\@ifundefined{KOMAClassName}{% if non-KOMA class
  \IfFileExists{parskip.sty}{%
    \usepackage{parskip}
  }{% else
    \setlength{\parindent}{0pt}
    \setlength{\parskip}{6pt plus 2pt minus 1pt}}
}{% if KOMA class
  \KOMAoptions{parskip=half}}
\makeatother
% Make \paragraph and \subparagraph free-standing
\makeatletter
\ifx\paragraph\undefined\else
  \let\oldparagraph\paragraph
  \renewcommand{\paragraph}{
    \@ifstar
      \xxxParagraphStar
      \xxxParagraphNoStar
  }
  \newcommand{\xxxParagraphStar}[1]{\oldparagraph*{#1}\mbox{}}
  \newcommand{\xxxParagraphNoStar}[1]{\oldparagraph{#1}\mbox{}}
\fi
\ifx\subparagraph\undefined\else
  \let\oldsubparagraph\subparagraph
  \renewcommand{\subparagraph}{
    \@ifstar
      \xxxSubParagraphStar
      \xxxSubParagraphNoStar
  }
  \newcommand{\xxxSubParagraphStar}[1]{\oldsubparagraph*{#1}\mbox{}}
  \newcommand{\xxxSubParagraphNoStar}[1]{\oldsubparagraph{#1}\mbox{}}
\fi
\makeatother


\usepackage{longtable,booktabs,array}
\newcounter{none} % for unnumbered tables
\usepackage{calc} % for calculating minipage widths
% Correct order of tables after \paragraph or \subparagraph
\usepackage{etoolbox}
\makeatletter
\patchcmd\longtable{\par}{\if@noskipsec\mbox{}\fi\par}{}{}
\makeatother
% Allow footnotes in longtable head/foot
\IfFileExists{footnotehyper.sty}{\usepackage{footnotehyper}}{\usepackage{footnote}}
\makesavenoteenv{longtable}
\usepackage{graphicx}
\makeatletter
\newsavebox\pandoc@box
\newcommand*\pandocbounded[1]{% scales image to fit in text height/width
  \sbox\pandoc@box{#1}%
  \Gscale@div\@tempa{\textheight}{\dimexpr\ht\pandoc@box+\dp\pandoc@box\relax}%
  \Gscale@div\@tempb{\linewidth}{\wd\pandoc@box}%
  \ifdim\@tempb\p@<\@tempa\p@\let\@tempa\@tempb\fi% select the smaller of both
  \ifdim\@tempa\p@<\p@\scalebox{\@tempa}{\usebox\pandoc@box}%
  \else\usebox{\pandoc@box}%
  \fi%
}
% Set default figure placement to htbp
\def\fps@figure{htbp}
\makeatother



\ifLuaTeX
\usepackage[bidi=basic,shorthands=off]{babel}
\else
\usepackage[bidi=default,shorthands=off]{babel}
\fi
\ifLuaTeX
  \usepackage{selnolig} % disable illegal ligatures
\fi


\setlength{\emergencystretch}{3em} % prevent overfull lines

\providecommand{\tightlist}{%
  \setlength{\itemsep}{0pt}\setlength{\parskip}{0pt}}



 


\KOMAoption{captions}{tableheading}
\makeatletter
\@ifpackageloaded{caption}{}{\usepackage{caption}}
\AtBeginDocument{%
\ifdefined\contentsname
  \renewcommand*\contentsname{Spis treści}
\else
  \newcommand\contentsname{Spis treści}
\fi
\ifdefined\listfigurename
  \renewcommand*\listfigurename{Spis rycin}
\else
  \newcommand\listfigurename{Spis rycin}
\fi
\ifdefined\listtablename
  \renewcommand*\listtablename{Spis tabel}
\else
  \newcommand\listtablename{Spis tabel}
\fi
\ifdefined\figurename
  \renewcommand*\figurename{Rysunek}
\else
  \newcommand\figurename{Rysunek}
\fi
\ifdefined\tablename
  \renewcommand*\tablename{Tabela}
\else
  \newcommand\tablename{Tabela}
\fi
}
\@ifpackageloaded{float}{}{\usepackage{float}}
\floatstyle{ruled}
\@ifundefined{c@chapter}{\newfloat{codelisting}{h}{lop}}{\newfloat{codelisting}{h}{lop}[chapter]}
\floatname{codelisting}{Wykaz}
\newcommand*\listoflistings{\listof{codelisting}{Spis wykazów}}
\makeatother
\makeatletter
\makeatother
\makeatletter
\@ifpackageloaded{caption}{}{\usepackage{caption}}
\@ifpackageloaded{subcaption}{}{\usepackage{subcaption}}
\makeatother
\usepackage{bookmark}
\IfFileExists{xurl.sty}{\usepackage{xurl}}{} % add URL line breaks if available
\urlstyle{same}
\hypersetup{
  pdftitle={Modelowanie Procesów Transportu w Środowisku},
  pdflang={pl},
  colorlinks=true,
  linkcolor={blue},
  filecolor={Maroon},
  citecolor={Blue},
  urlcolor={Blue},
  pdfcreator={LaTeX via pandoc}}


\title{Modelowanie Procesów Transportu w Środowisku}
\usepackage{etoolbox}
\makeatletter
\providecommand{\subtitle}[1]{% add subtitle to \maketitle
  \apptocmd{\@title}{\par {\large #1 \par}}{}{}
}
\makeatother
\subtitle{Ćwiczenia laboratoryjne}
\author{}
\date{}
\begin{document}
\maketitle


\section{Ćwiczenie 1. Od błądzenia przypadkowego do transportu
zanieczyszczeń w
środowisku}\label{ux107wiczenie-1.-od-bux142ux105dzenia-przypadkowego-do-transportu-zanieczyszczeux144-w-ux15brodowisku}

Celem ćwiczenia jest opanowanie podstaw matematycznego opisu procesów
transportu materii w środowisku oraz zrozumienie założeń, interpretacji
i zastosowań wybranych modeli poprzez analizę i rozwiązywanie
przykładowych problemów.

Ćwiczenie polega na symulacji ruchu pojedynczych cząstek lub ich
zespołów poprzez realizację ciągów skoków wykonywanych w stałych
(\(\Delta t\)) lub zmiennych (\(\Delta t_n\)) odstępach czasu. Przyrosty
przemieszczenia oznaczamy jako \(\Delta x_n\) i losujemy z zadanych
rozkładów prawdopodobieństwa.

Położenie cząstki (\(x_{n+1}\)) w kolejnych chwilach czasu (\(t_{n+1}\))
wyznaczamy jako: \[ t_{n+1}=t_n+\Delta t_n,\]
\[x_{n+1}=x_n+\Delta x_n.\]

Dla stałych przyrostów czasu mamy: \(t_n=n\Delta t.\)

Wszystkie cząstki rozpoczynają ruch w punkcie: \(x_0=0, y_0=0\).
Symulacje prowadzimy w \(M\) krokach czasowych dla \(N\) cząstek.

\subsection{\texorpdfstring{Błądzenie przypadkowe (\emph{Random Walk})
dla jednej
cząstki}{Błądzenie przypadkowe (Random Walk) dla jednej cząstki}}\label{bux142ux105dzenie-przypadkowe-random-walk-dla-jednej-czux105stki}

\subsubsection{Zadanie: symulacja ruchu w jednym
wymiarze}\label{zadanie-symulacja-ruchu-w-jednym-wymiarze}

Przeprowadź symulację ruchu jednej cząstki dla regularnych przyrostów
czasu \(\Delta t\) i przyrostów położenia losowanych z następujących
rozkładów.

\begin{itemize}
\item
  \textbf{Rozkład Bernoulliego} \[
  \Delta x_n=
  \begin{cases}
  -l, & p=\frac{1}{2},\\
  +l, & p=\frac{1}{2}.
  \end{cases}
  \]
\item
  \textbf{Rozkład jednostajny} \[\Delta x_n\sim U(-l,l)\]
\item
  \textbf{Rozkład normalny} \[\Delta x_n\sim\mathcal{N}(0,\sigma^2)
  \]
\item
  Opcjonalnie przeprowadź symulację dla innych rozkładów o średniej zero
  i skończonej wariancji, np. rozkładu Laplace'a.
\end{itemize}

Dla każdego przypadku: - Wygeneruj trajektorię cząstki. - Przedstaw
położenie cząstki w funkcji czasu. - Porównaj otrzymane trajektorie. -
Zinterpretuj wpływ rozkładu przyrostów na ruch cząstki.

\subsubsection{Zadanie: Błądzenie przypadkowe w dwóch
wymiarach}\label{zadanie-bux142ux105dzenie-przypadkowe-w-dwuxf3ch-wymiarach}

Powtórz symulacje w dwóch wymiarach przestrzennych. Zwizualizuj i
porównaj trajektorie na płaszczyźnie \((x,y)\).

\subsubsection{Zadanie: Rozkład końcowych położeń
cząstek}\label{zadanie-rozkux142ad-koux144cowych-poux142oux17ceux144-czux105stek}

Powtórz symulacje w jednym wymiarze dla \(M\in\{10^2,10^3,10^4\}\) oraz
\(N\in\{10^2,10^4\}\). Narysuj histogramy końcowych położeń dla
wszystkich sześciu kombinacji parametrów \(M\) i \(N\). Sprawdź zgodność
otrzymanych rozkładów z rozkładem normalnym.

\subsubsection{Zadanie: Średnie położenie i wariancja populacji
cząstek}\label{zadanie-ux15brednie-poux142oux17cenie-i-wariancja-populacji-czux105stek}

Dla rozkładów przyrostów wyznacz średnie położenie po
\(n\in\{1,\ldots,1000\}\) krokach czasowych: \[
\overline{x}(n)=\frac{1}{N}\sum_{i=1}^{N}x_n^i,
\] oraz wariancję: \[
\sigma_x^2(n)=
\frac{1}{N}\sum_{i=1}^{N}
\left(x_n^i-\overline{x}(n)\right)^2.
\]

\section{Ćwiczenie 2. Od błądzenia przypadkowego do adwekcji-dyfuzji
zanieczyszczeń}\label{ux107wiczenie-2.-od-bux142ux105dzenia-przypadkowego-do-adwekcji-dyfuzji-zanieczyszczeux144}

Celem ćwiczenia jest zrozumienie mechanizmu adwekcji i jego wpływu na
transport oraz rozprzestrzenianie się substancji w środowisku.

Student powinien umieć: - Rozróżniać adwekcję i dyfuzję. - Rozszerzyć
algorytm błądzenia przypadkowego o średnie przemieszczenie związane z
prędkością przepływu. - Analizować wpływ prędkości adwekcji \(u\) i
współczynnika dyfuzji \(D\). - Porównywać symulację cząstkową z
rozwiązaniem równania adwekcji-dyfuzji. - Interpretować ruch środka masy
chmury oraz zmianę wariancji.

Wprowadzenie adwekcji polega na dodaniu składnika deterministycznego: \[
{x_{n+1}=x_n+u\Delta t+\xi_n}
\] gdzie: - \(x_n\) --- położenie cząstki w chwili \(t_n\); - \(u\) ---
stała prędkość adwekcji; - \(\Delta t\) --- krok czasowy; - \(\xi_n\)
--- losowy przyrost położenia związany z dyfuzją.

Przyrosty losowe są niezależne i spełniają warunki:
\[\mathbb{E}[\xi_n]=0,\] \[
\operatorname{Var}(\xi_n)=2D\Delta t.
\] Przyjmujemy rozkład normalny: \[
\xi_n\sim\mathcal{N}(0,2D\Delta t).
\]

Adwekcja przesuwa średnie położenie chmury cząstek, natomiast dyfuzja
zwiększa rozrzut ich położeń.

\subsubsection{Równanie
adwekcji-dyfuzji}\label{ruxf3wnanie-adwekcji-dyfuzji}

Makroskopowy opis transportu przy stałej prędkości przepływu i stałym
współczynniku dyfuzji przedstawia równanie: \[
\frac{\partial C}{\partial t} + u\frac{\partial C}{\partial x} =D\frac{\partial^2 C}{\partial x^2},
\] gdzie: - \(C(x,t)\) --- stężenie substancji, - \(u\) --- prędkość
adwekcji, - \(D\) --- współczynnik dyfuzji.

\subsubsection{Zadanie: Wprowadzenie adwekcji do modelu
cząstkowego}\label{zadanie-wprowadzenie-adwekcji-do-modelu-czux105stkowego}

Rozważ transport cząstek rozpoczynających ruch w punkcie \(x_0=0\).
Wykonaj symulacje dla, np.: \(N=10^4,M=1000.\) Przykładowe parametry do
symulacji:

{\def\LTcaptype{none} % do not increment counter
\begin{longtable}[]{@{}lr@{}}
\toprule\noalign{}
Parametr & Wartość \\
\midrule\noalign{}
\endhead
\bottomrule\noalign{}
\endlastfoot
Krok czasowy \(\Delta t\) & \(1\ \mathrm{s}\) \\
Współczynnik dyfuzji \(D\) & \(0{,}5\ \mathrm{m^2/s}\) \\
Prędkość adwekcji \(u\) & \(0;\ 0{,}1;\ 0{,}5;\ 1\ \mathrm{m/s}\) \\
\end{longtable}
}

Dla każdej wartości \(u\): 1. Wygeneruj trajektorie cząstek. 2. Narysuj
położenie wybranych cząstek w funkcji czasu. 3. Porównaj rozkłady
położeń dla różnych chwil czasu.

\subsubsection{Zadanie: Średnie położenie i wariancja populacji
cząstek}\label{zadanie-ux15brednie-poux142oux17cenie-i-wariancja-populacji-czux105stek-1}

Dla każdego kroku czasowego wyznacz średnie położenie oraz wariancję:

\[
\mathbb{E}[x(t)]=ut,
\]

\[
\operatorname{Var}[x(t)]=2Dt.
\] Porównaj z wartościami teoretycznymi, zbadaj liniowość otrzymanych
rezultatów.

\subsubsection{Zadanie: Rozkład położeń cząstek a rozwiązanie
analityczne}\label{zadanie-rozkux142ad-poux142oux17ceux144-czux105stek-a-rozwiux105zanie-analityczne}

Porównaj wynik symulacji z rozwiązaniem analitycznym Dla początkowego
rozkładu punktowego \(C(x,0)=\delta(x)\) rozwiązanie równania
adwekcji-dyfuzji wynosi:

\[
C(x,t)= \frac{1}{\sqrt{4\pi Dt}}  \exp\left(-\frac{(x-ut)^2}{4Dt}\right).
\]

\begin{center}\rule{0.5\linewidth}{0.5pt}\end{center}

\section{Ćwiczenie 3. Analiza wyników testu znacznikowego w
rzece}\label{ux107wiczenie-3.-analiza-wynikuxf3w-testu-znacznikowego-w-rzece}

Do rzeki wlano wodny roztwór \(m = 10\) kg KBr. Próbki wody pobierano z
rzeki w odległości \(x = 2000 m\) poniżej miejsca iniekcji
(wprowadzenia) znacznika. Załączone dane przedstawiają stężenia
znacznika (jonu Br-) \(C_i\) dla kolejnych próbek wody pobranych w
chwilach \(t_i\).

\subsubsection{Zadanie: wykreślenie krzywej
przejścia}\label{zadanie-wykreux15blenie-krzywej-przejux15bcia}

Należy narysować krzywą przejścia znacznika (BTC -- breakthrough curve),
czyli zależność \(C_i = C(t_i)\), a następnie wyliczyć momenty tej
krzywej, \(M_k\), dla \(k=0, 1, 2\). Do całkowania określonej dyskretnie
funkcji \(C(t_i)\) zastosować metodę trapezów.

\[
M_k = \int_0^\infty t^k C(t)\,dt.
\]

\subsubsection{Zadanie: wyznaczenie wielkości charakteryzujących krzywą
przejścia}\label{zadanie-wyznaczenie-wielkoux15bci-charakteryzujux105cych-krzywux105-przejux15bcia}

\begin{itemize}
\item
  Natężenie przepływu wody w rzece (Q) {[}m³/s{]} \[
  Q=\frac{m_{\mathrm{Br}^-}}{M_0}
  \]
\item
  Średni czas przejścia znacznika \((\bar{t})\) {[}s{]} \[
  \bar{t}=\frac{M_1}{M_0}
  \]
\item
  Wariancja czasu przejścia znacznika \((\sigma_t^2)\) {[}s²{]}
  \[\sigma_t^2=\frac{M_2}{M_0}-\left(\frac{M_1}{M_0}\right)^2\]
\end{itemize}

\subsubsection{Zadanie: określenie wpływu niepewności analitycznych na
niepewność charakterystyk krzywej
przejścia}\label{zadanie-okreux15blenie-wpux142ywu-niepewnoux15bci-analitycznych-na-niepewnoux15bux107-charakterystyk-krzywej-przejux15bcia}

Stężenia jonu bromkowego określono ze względną niepewnością 5\%. Celem
tej części ćwiczenia jest sprawdzenie, jak te niepewności propagują się
do charakterystyk wyznaczonych w poprzedniej części. Stosujemy do tego
metodę Monte Carlo, generując \[N=10^4\] krzywych przejścia:

\[
C_i^j \sim \mathcal{N}\left(C_i,(0{,}05C_i)^2\right).
\]

Dla każdej z \(N\) otrzymanych w ten sposób krzywych przejścia
wyznaczamy \(Q\), \(\bar{t}\), \(\sigma_t^2\), a następnie dla
otrzymanych rozkładów tych wielkości wyznaczamy wartości średnie i
odchylenia standardowe, np. dla \(Q\):

\textbf{Wartość średnia:} \(\mu_Q=\frac{1}{N}\sum_{j=1}^{N}Q^j\)

\textbf{Odchylenie standardowe:}
\(s_Q=\sqrt{\frac{1}{N-1}\sum_{j=1}^{N}(Q^j-\mu_Q)^2}\)

\begin{center}\rule{0.5\linewidth}{0.5pt}\end{center}

\section{Ćwiczenie 4. Modelowanie transportu zanieczyszczenia w rzece
(Crank--Nicholson lub
QUICKEST)}\label{ux107wiczenie-4.-modelowanie-transportu-zanieczyszczenia-w-rzece-cranknicholson-lub-quickest}

Celem ćwiczenia jest numeryczne modelowanie transportu znacznika w rzece
z wykorzystaniem równania adwekcyjno-dyspersyjnego. Student powinien
umieć zaimplementować wybraną metodę numeryczną, ocenić wpływ
dyskretyzacji na otrzymane rozwiązanie oraz porównać wyniki symulacji z
danymi pomiarowymi.

\subsubsection{Równanie adwekcyjno-dyspersyjne
(RAD)}\label{ruxf3wnanie-adwekcyjno-dyspersyjne-rad}

\[
\frac{\partial c}{\partial t}  +U\frac{\partial c}{\partial x}  -D\frac{\partial^2 c}{\partial x^2}=0
\] gdzie:

\begin{itemize}
\tightlist
\item
  \(c\) -- stężenie zanieczyszczenia w wodzie {[}kg/m³{]},
\item
  \(U\) -- prędkość wody {[}m/s{]},
\item
  \(D\) -- współczynnik dyspersji {[}m²/s{]}.
\end{itemize}

\begin{enumerate}
\def\labelenumi{\arabic{enumi}.}
\item
  Warunki brzegowe

  \begin{itemize}
  \tightlist
  \item
    \textbf{Górny brzeg obszaru symulacji} --- warunek Dirichleta
    \(c(0,t)=0\)
  \item
    \textbf{Dolny brzeg obszaru symulacji} --- warunek Dirichleta
    \(c(L,t)=0\)
  \end{itemize}

  Znacznik nie dociera do dolnego brzegu symulowanego obszaru, ale nie
  możemy z góry założyć, że warunek będzie spełniony dla wybranego
  \(L\). Stosujemy więc zabieg polegający na ustawieniu dolnego brzegu
  na tyle daleko, że zanieczyszczenie nie osiągnie go w trakcie trwania
  symulacji. Wielkość \(L\) ustalamy więc dynamicznie, chociaż może się
  zdarzyć, że wybrana początkowo wartość będzie odpowiednio duża i nie
  zajdzie potrzeba rozszerzania obszaru symulacji.
\item
  Warunki początkowe \[
   c(0,j)=c_{\mathrm{in}}
   \quad \text{dla} \quad j=\frac{l_{\mathrm{in}}}{\Delta x}+1 ,
  \]

  \[
   c(0,j)=0 \quad \text{dla} \quad
   j\neq\frac{l_{\mathrm{in}}}{\Delta x}+1
   \]

  Stężenie początkowe wynosi
  \(c_{\mathrm{in}}=\frac{m}{A\cdot\Delta x}\), gdzie:

  \begin{itemize}
  \tightlist
  \item
    \(m\) --- całkowita masa zanieczyszczenia,
  \item
    \(A\) --- pole przekroju poprzecznego kanału,
  \item
    \(\Delta x\) --- krok przestrzenny.
  \end{itemize}
\item
  Wartości parametrów
\end{enumerate}

{\def\LTcaptype{none} % do not increment counter
\begin{longtable}[]{@{}lrl@{}}
\toprule\noalign{}
Parametr & Wartość & Opis \\
\midrule\noalign{}
\endhead
\bottomrule\noalign{}
\endlastfoot
\(L_{\mathrm{in}}\) & 10 m & Odległość punktu iniekcji \\
\(L_{\mathrm{out}}\) & 2100 m & Długość obszaru za punktem iniekcji \\
\(A\) & 10 m² & Pole przekroju poprzecznego kanału \\
\(U\) & 0,1 m/s & Prędkość wody \\
\(D\) & 1 m²/s & Współczynnik dyspersji \\
\(m\) & 10 kg & Całkowita masa zanieczyszczenia \\
\end{longtable}
}

Całkowita długość obszaru obliczeniowego wynosi:

\[
L=L_{\mathrm{in}}+L_{\mathrm{out}}=2110\ \mathrm{m}.
\]

\subsubsection{Rozwiązanie RAD metodą
QUICKEST}\label{rozwiux105zanie-rad-metodux105-quickest}

W metodzie QUICKEST stężenie w kolejnym kroku czasowym obliczamy ze
wzoru: \[
\begin{aligned}
    c_j^{n+1}=c_j^n &+\left[C_d(1-C_a)-\frac{C_a}{6} (C_a^2-3C_a+2)\right]c_{j+1}^n\\
                    &-\left[C_d(2-3C_a)-\frac{C_a}{2}(C_a^2-2C_a-1)\right]c_j^n\\
                    &+\left[C_d(1-3C_a)-\frac{C_a}{2}(C_a^2-C_a-2)\right]c_{j-1}^n\\
                    &+\left[C_dC_a+\frac{C_a}{6}(C_a^2-1)\right]c_{j-2}^n
\end{aligned}
\]

gdzie:

\[
C_a=\frac{U\Delta t}{\Delta x}: \quad \text{ adwekcyjna liczba Couranta,}
\] \[
C_d=\frac{D\Delta t}{(\Delta x)^2}: \quad \text{ dyfuzyjna liczba Couranta.}
\]

\subsubsection{Zadanie: porównanie symulacji z danymi
pomiarowymi}\label{zadanie-poruxf3wnanie-symulacji-z-danymi-pomiarowymi}

\begin{enumerate}
\def\labelenumi{\arabic{enumi}.}
\item
  Zaimplementuj metodę QUICKEST i przeprowadź symulację transportu
  znacznika w rzece dla podanych parametrów.
\item
  Dobierz krok przestrzenny \(\Delta x\) oraz krok czasowy \(\Delta t\).
  Sprawdź warunek Couranta--Friedrichsa--Lewy'ego (CFL) oraz
  ograniczenia stabilności zastosowanego schematu numerycznego.
\item
  Wykonaj symulacje dla wybranych wartości \(\Delta t\), zarówno
  spełniających, jak i naruszających przyjęte warunki stabilności.
  Porównaj profile stężenia i oceń stabilność wyników.
\item
  Wyznacz krzywą przejścia znacznika w punkcie pomiarowym \(x_n\). Punkt
  pomiarowy znajduje się 2000 m poniżej miejsca iniekcji, czyli:

  \[
  x_n=L_{\mathrm{in}}+2000\ \mathrm{m}=2010\ \mathrm{m}.
  \]
\item
  Porównaj symulowaną krzywą przejścia z danymi pomiarowymi. W razie
  potrzeby interpoluj stężenia symulowane do chwil wykonania pomiarów.
\item
  Oblicz błędy symulacji, wykorzystując średni błąd bezwzględny (MAE)
  oraz pierwiastek z błędu średniokwadratowego (RMSE).
\item
  Przedstaw na wykresie krzywe przejścia uzyskane z symulacji i
  pomiarów.
\item
  Sformułuj wnioski dotyczące wpływu kroku przestrzennego i czasowego na
  stabilność oraz dokładność rozwiązania.
\end{enumerate}

\subsubsection{Zadanie: analiza zbieżności
rozwiązania.}\label{zadanie-analiza-zbieux17cnoux15bci-rozwiux105zania.}

Wykonaj symulacje metodą QUICKEST dla kolejno zagęszczanych siatek
przestrzennych, np. \(\Delta x \in \{20, 10, 5, 2{,}5\}\ \mathrm{m}\),
odpowiednio dobierając krok czasowy \(\Delta t\) i kontrolując
stabilność numeryczną. Porównaj otrzymane krzywe przejścia znacznika,
oblicz błędy RMSE względem rozwiązania na najgęstszej siatce i przedstaw
zależność błędu od \(\Delta x\). Oceń zbieżność rozwiązania oraz wpływ
zagęszczania siatki na dokładność i koszt obliczeń.




\end{document}
